% Experimental setup (journal format, 2026-11-19). Numbers from paper/tables (scripts/build_paper.py).
\section{Experimental setup}
\label{sec:experiments}
\label{sec:setup}

\subsection{Benchmark instances}
\emph{Suite A} reproduces the instance sizes of \citet{jin2026imoeadqn}: 3, 6 or 10 edge servers, 100, 300 or 500 tasks,
and 20--250 devices that each submit a layered random DAG. Workloads are 1--5 Gcycles and data 500--3000 KB as in Jin et
al., used as given and multiplied by 20 to make communication matter, which gives 12 test instances; 36 further
instances generated with other seeds were used only to set the method's parameters. \emph{Suite B} takes five
scientific workflow families from WfCommons \citep{coleman2022wfcommons}, namely Cybershake, Epigenomics, Inspiral,
Montage and Sipht, with 100, 400 and 1000 tasks, submitted by one device in Jin et al.'s small scenario (3 edge servers,
6 devices). Workloads and data sizes are mapped into Jin et al.'s ranges by an order-preserving logarithmic map, every
quantity has a relative interval half-width of 5--30\%, and the data sizes are multiplied by 1, 20 or 100 to vary the
communication-to-computation ratio, giving 45 instances. Fifteen further workflows with 2000 tasks test scalability.
Table~\ref{tab:params_main} lists the main parameters.

% generated by scripts/build_paper.py; do not edit by hand
\begin{table}[!htbp]
\centering
\footnotesize
\caption{Main parameters (Jin et al.'s small scenario). All environment values and the remaining settings are in Table~\Sref{tab:params} of the supplementary material.}
\label{tab:params_main}
\setlength{\tabcolsep}{3pt}
\begin{tabular}{l>{\raggedright\arraybackslash}p{0.52\columnwidth}}
\toprule
Parameter & Value \\
\midrule
\multicolumn{2}{l}{\textit{Environment}} \\
Device / edge CPU & $0.45$--$1.18$ / $78.5$--$150$ GHz \\
Cloud CPU per task & $[18, 22]$ GHz \\
Bandwidth / edge--cloud & $[9, 12]$ MHz / $[408, 624]$ Mbps \\
Prices $\eta_E$ / $\eta_C$ & 0.001 \$/(GHz\,s) / 0.01 \$/s \\
Workload / data & 1--5 Gcycles / 500--3000 KB \\
\midrule
\multicolumn{2}{l}{\textit{\ours{}}} \\
$N$ / $G_{\max}$ & 105 / 1000 \\
$(\alpha, \beta, \gamma)$ & $(5/6, 1/6, 0)$ \\
Weight vectors & 105 (Das--Dennis, $q = 13$) \\
ICD neighbours $k$ & 14 \\
Mating $T$ / $p_{local}$ & 10 / 1 \\
$p_c$ / mutation & 0.7, 0.9 / $1/K$ \\
\midrule
\multicolumn{2}{l}{\textit{Protocol}} \\
Runs & 30 (\ours{}, IMOEA/D, IMOEA-DQN (fixed)), 10 others \\
Evaluations per run & 105,105 \\
\bottomrule
\end{tabular}
\end{table}


\subsection{Compared algorithms}
We compare \ours{} with eight methods, all implemented on the same encoding, variation operators and evaluator and given
the same population ($N = 105$) and budget of $N(G_{\max}+1)$ evaluations with $G_{\max} = 1000$:
\begin{itemize}
\item \emph{IMOEA/D}, the interval decomposition MOEA with weight-vector adaptation, with the original settings reported
  by \citet{cui2024imeaied} ($T = 10$, adaptation from $G_{\max}/2$ every $G_{\max}/6$);
\item \emph{IMOEA/D-AWN} \citep{jin2023imoeadawn}, which adapts weight vectors and neighbourhood sizes by interval sparsity;
\item \emph{IMEA-IED} \citep{cui2024imeaied}, with information-entropy interval dominance and niche selection;
\item \emph{IP-MOEA} \citep{limbourg2005ipmoea}, with interval Pareto dominance and hypervolume-contribution crowding;
\item \emph{NSGA-II} \citep{deb2002nsga2} on interval midpoints (mid) or with interval dominance (int), and
  \emph{NSGA-III} \citep{deb2014nsga3} on midpoints;
\item \emph{IMOEA-DQN} \citep{jin2026imoeadqn} as published, with an online DQN choosing $(\alpha, \beta, \gamma)$ in every
  generation; and its \emph{Jin backbone} run with the fixed action $(5/6, 1/6, 0)$, which isolates the effect of
  Rank C and the mating.
\end{itemize}
Six list-scheduling heuristics are reported in Table~\Sref{tab:heuristics} of the supplementary material: all tasks
local, on the nearest edge server or in the cloud, and three HEFT-style schedulers that place each task by earliest
finish time, lowest energy or lowest cost \citep{topcuoglu2002heft}.

\subsection{Performance metrics and statistics}
Each instance is normalised once, over the final fronts of all compared methods, so that every objective lies in
$[0, 1]$. With $R$ the non-dominated union of the midpoints of all final fronts, the inverted generational distance of
a front $P$ is
\begin{equation}
\mathrm{IGD}(P) = \frac{1}{|R|} \sum_{r \in R} \min_{x \in P} \lVert r - m(\tilde F(x)) \rVert_2,
\label{eq:igd}
\end{equation}
and, as in \citet{jin2026imoeadqn}, the hypervolume is the mean of the hypervolumes of the lower and the upper bounds,
$\mathrm{HV}(P) = \tfrac{1}{2}\big(\mathrm{HV}(\{F^L(x)\}) + \mathrm{HV}(\{F^U(x)\})\big)$, with reference point $1.1$ in
every objective. IGD measures convergence and coverage (lower is better), HV the dominated volume (higher is better).
The imprecision $\sum_{x \in P} \sum_m w(\tilde F_m(x))$ measures the width of the predicted intervals.

\ours{}, IMOEA/D and the Jin backbone were run 30 times per instance and the other methods 10 times. Each method is
compared with \ours{} on each instance by the Wilcoxon rank-sum test with Holm's correction \citep{holm1979} over the
compared methods at $\alpha = 0.05$, and all methods are compared across instances by their mean Friedman ranks with the
Nemenyi critical difference \citep{demsar2006}.
